\vspace{-2mm}
\section{Related Work}
\label{sec:related}

\noindent \textbf{Parametric optimization:} For some flash loan attack cases,
researchers~\cite{cao2021flashot, qin2021attacking} manually extracted math formulas
of function candidates, manually defined related parameter constraints, and 
used an off-the-shelf optimizer to search for parameters which yield the best profit. 
However, this technique requires
significant manual efforts and expert knowledge of the underlying DeFi
protocols. Consequently, it becomes impractical for checking a large number of potential attack vectors. 
Note that our benchmark set contains significantly more flash loan attacks than 
prior work~\cite{qin2021attacking,cao2021flashot}, i.e., $18$ versus $2$ in~\cite{qin2021attacking} and $9$ in~\cite{cao2021flashot}. 

\noindent \textbf{Static Analysis:} 
Slither~\cite{feist2019slither}, Securify~\cite{tsankov2018securify},
Zeus~\cite{kalra2018zeus}, Park~\cite{zheng2022park} and
SmartCheck~\cite{tikhomirov2018smartcheck} apply static analysis techniques to
verify smart contracts. 
%They detect pre-defined patterns of vulnerabilities and
%over-approximate program states, which inevitably cause false positives and
%negatives. 
There are also several works that use
symbolic execution~\cite{king1976symbolic} to explore the program states of a
smart contract, looking for an execution path that violates a user-defined
invariant, e.g., Mythril~\cite{mythril}, Oyente~\cite{luu2016making},
FairCon~\cite{liu2020towards}, ETHBMC~\cite{frank2020ethbmc},
SmartCopy~\cite{feng2019precise}, and Manticore~\cite{mossberg2019manticore}.
These techniques tend to operate with one contract at a time and therefore cannot handle flash loan attacks that involve multiple contracts.
These techniques also may suffer from the complicated logics of the DeFi contracts, and cause
path explosion.
{\name} uses its novel synthesis-via-approximation techniques to avoid these issues. 


\noindent \textbf{Fuzzing:}
ContractFuzzer~\cite{jiang2018contractfuzzer}, sFuzz~\cite{nguyen2020sfuzz},
ContraMaster~\cite{wang2020oracle}, SMARTIAN~\cite{choi2021smartian} and ItyFuzz~\cite{shou2023ityfuzz}
introduce novel fuzzing techniques to discover vulnerabilities in smart
contracts. However, these techniques either only work on one contract or focus
on specific vulnerabilities like re-entrancy. Moreover, flash loan attack vectors can contain 
up to $8$ actions and $7$ integer parameters, which is unlikely to be found by random fuzzing. 


\noindent \textbf{Flash Loan Attacks:} 
Prior works~\cite{qin2021attacking,cao2021flashot,werapun2022flash} study specific flash loan attacks and 
manually analyze the attack vectors. Some tools~\cite{xia2023detecting, ramezany2023midnight, 
wang2022defiscanner} 
are designed to monitor flash loan attacks after they happened. Some researchers investigate other usage of 
flash loans including arbitrage~\cite{wang2021towards} or wash trading~\cite{,gan2022understanding}. To the best
of our knowledge, {\name} is the first tool that automatically synthesizes complicated flash loan attack traces.
It shows successes in real-world DeFi protocols under audit. 
